% ════════════════════════════════════════════════════════════════════════════
\section{Related work}
\label{sec:related}
% ════════════════════════════════════════════════════════════════════════════

\textbf{Coercion resistance.} The term is due to Juels, Catalano and
Jakobsson~\cite{juels-coercion}, in electronic voting: a scheme is
coercion-resistant if a voter \emph{cannot prove} how they voted, so the coercer
cannot verify compliance and the reason to coerce collapses. Our sense is
complementary. There the protected act is a \emph{choice} and safety comes from
\emph{unprovability}; here the protected object is a \emph{bearer secret}, and the
danger is not unverifiable compliance but that a released holder remains a racer
worth killing. We inherit the word, not the model.

\textbf{Secrets recognized but not articulated.} The tacit premise
(Assumption~\ref{ass:tkba}) has direct antecedents. Bojinov, Sanchez, Reber, Boneh
and Lincoln~\cite{bojinov-rubberhose} train an authentication secret by
\emph{implicit learning} so that even a cooperating, coerced user cannot
articulate it --- the closest prior instantiation of our TKBA, with TGPO carrying
the idea from an authentication token to a self-custody key behind a memory-hard
time-gate. The broader phenomenon is old: keystroke dynamics treat typing rhythm
as a biometric~\cite{monrose-keystroke}, and WWII traffic analysis identified
wireless operators by the ``fist'' of their Morse keying~\cite{kahn-codebreakers}.

\textbf{Deniability versus unprovability.} Deniable
encryption~\cite{canetti-deniable} and the duress-PIN schemes it inspired let a
coerced user \emph{plausibly lie}. NNPP is the sharp negative counterpart: rather
than proving a false positive (``this is the key''), the holder would need to prove
a true negative (``no other key exists''), which cannot be done. Deniability offers
a lie the coercer cannot disprove; we show the holder cannot furnish the assurance
the coercer needs, which is why the coercer's only remaining lever is the victim.
Such schemes also rest on the adversary's ignorance and so fail Kerckhoffs
(\S\ref{sub:selfdestruct}).

\textbf{Time-gates.} The construction's wall-clock gate is a time-lock puzzle in
the sense of Rivest, Shamir and Wagner~\cite{rsw1996}; verifiable delay
functions~\cite{boneh-vdf} are the modern publicly checkable form, and either may
instantiate the outsourced time-lock the orbit optionally uses.

\textbf{Kidnap economics and marginal deterrence.} The Lemma's shape is not
new, and it would be wrong to present it as though it were. Selten's kidnapping
game~\cite{selten-kidnapping} is the canonical formalization of the encounter we
model, and the criminological literature on \emph{marginal deterrence} ---
Stigler~\cite{stigler-enforcement}, formalized by
Shavell~\cite{shavell-marginal-deterrence} --- has long carried the observation
that an offender facing no incremental sanction has no reason to stop short of the
worse act, with \emph{witness elimination} as its standard illustration. At that
level of abstraction the Deadly-Race Lemma is a witness-elimination argument in
which what the victim retains is not testimony but a claim on a rivalrous prize.
We claim no priority over either literature.

What is worth stating precisely is \emph{where} our setting departs from Selten's,
because the departure is why our conclusion differs. Selten assumes the
kidnapper's threat to execute is credible \emph{even though executing cannot
improve his position} --- correct in classic kidnap-for-ransom, where the prize is
transferred by a third party and is not contestable once paid, so the hostage's
survival threatens nothing the kidnapper holds. Bitcoin self-custody removes
exactly that property: the prize is rivalrous and bearer, so a design leaving
$\Victim$ a feasible path leaves it \emph{contestable after release}. In Selten's
terms our $(1-p_A)$ is identically zero; the Deadly Race is his game with a
non-degenerate post-release contest, and that single structural change turns
execution from a threat instrument into an action with positive expected value.
The contribution is therefore not the mechanism but its locus: marginal deterrence
studies how a \emph{legislator's} penalty schedule moves this margin, whereas a
\emph{custody designer} moves the same margin --- and a widely shipped class of
designs moves it the wrong way.

Two qualifications. Iqbal, Masson and Abbott~\cite{iqbal-selten-extension} extend
Selten by letting the capture probability rise when the hostage is executed, which
is our $c$ made explicit; our claim is that the gain side can exceed $c$ for large
$W$, not that $c$ is small, which is why Lemma~\ref{lem:drl} is a threshold
condition. And the deterrence literature in this vein --- Lapan and Sandler on
precommitment to non-negotiation~\cite{lapan-sandler-bargain} --- addresses
whether the encounter \emph{occurs}, whereas our results concern the game once it
has (\S\ref{sec:classify}).
